\documentclass[11pt]{article}
\usepackage{amsmath}
\usepackage{amssymb}
\usepackage[dvipsnames]{xcolor}
\usepackage{graphicx}
\usepackage{newpxtext}
\usepackage{eulervm}
\DeclareMathOperator{\diag}{\mathrm{diag}}
\DeclareMathOperator{\cov}{\texttt cov}
\long\def\ans#1{{\color{Blue}{\em #1}}}
\usepackage[margin=1in]{geometry}
\title{Homework 5}
\author{Your name here}
\date{\today{}}

\begin{document}
\maketitle

\section{MLE for independent samples from the Bernoulli distribution}

\emph{(Adapted from FCML Ex.~2.9)}

\subsection*{Context \& purpose}

\noindent The Bernoulli MLE---``the estimate is just the sample
proportion''---is the simplest instance of a pattern you will repeat all term:
write a likelihood, maximize its logarithm, and read off the estimator. Doing it
carefully once here makes the Gaussian, Poisson, and logistic cases mechanical.

\subsection*{Learning objectives}

\noindent By the end of this problem you can:
\begin{itemize}
    \item \textbf{O1} Write the likelihood of an i.i.d.\ dataset from its
        per-sample mass function. \emph{(Bloom: Apply --- Correctness step A)}
    \item \textbf{O2} Show a log-likelihood has a unique maximum. \emph{(Bloom:
        Analyze --- Correctness step B)}
    \item \textbf{O3} Derive the maximum-likelihood estimator in closed form.
        \emph{(Bloom: Analyze --- Correctness step C)}
    \item \textbf{O4} Interpret the estimator---why it equals the sample mean,
        when that holds, and when it misleads. \emph{(Bloom: Evaluate ---
        Interpretation rubric)}
\end{itemize}

\subsection*{The problem}

\noindent The Bernoulli distribution can be used to model events with binary outcomes
(e.g., coin tosses). For a random variable $X$ that can take on two values: 0 or
1, the probability mass function describing the outcome (i.e., the likelihood) is given by
\begin{align}
P(x|r) = r^x(1-r)^{1-x}
\end{align}
where $r$ is the probability of $X$ taking on the value 1.

\begin{enumerate}
    \item Suppose you observe $N$ values $x_1, \dots, x_N$ sampled from the
        Bernoulli distribution. Write down the expression for the likelihood
        $L$ of this dataset, assuming that the samples are independent
        conditioned on $r$.
    \item Show that $L$ has a unique maximum.
    \item Now that you have shown that $L$ has a unique maximum, derive an
        expression for the maximum likelihood estimator $\hat{r}$
        for the parameter $r$ in the likelihood $L$.
\end{enumerate}

\subsection*{What you may and may not use}

\noindent \textbf{Any mathematically valid derivation that reaches the required
result earns full credit}---there is no privileged path. You may maximize $L$
directly or maximize $\log L$; you may argue uniqueness by the second-derivative
(concavity) test, by a strict-concavity argument, or by showing the score
function has a single sign change on $(0,1)$. You must \emph{show the steps that
justify each move}: a bare final answer earns credit only for the result step.

\subsection*{How you'll be assessed (criteria shown up front)}

\noindent Grading follows the shared analytic rubric in \texttt{rubric.md}
(shipped with this assignment), not the old complete/incomplete gate:
\begin{itemize}
    \item \textbf{Correctness (60\%)} --- a per-step derivation rubric (setup
        $\to$ uniqueness $\to$ result); each step earns partial credit on the
        \emph{validity of the move}, and any sound alternative route to the same
        result earns full credit.
    \item \textbf{Interpretation (30\%)} --- the paragraph below, scored
        Claim / Evidence / Reasoning / Limits (0--3 each).
    \item \textbf{Process (10\%)} --- clear notation, stated assumptions, and
        legible justification.
\end{itemize}
Under the specifications-grading bundle, ``Derivation = PASS'' requires steps
A--C each at Proficient or above; otherwise you get one revise-and-resubmit.

\subsection*{Required interpretation}

\noindent After deriving $\hat{r}$, write \textbf{one short paragraph} that:
(a) explains \emph{why} $\hat{r} = \frac{1}{N}\sum_{i} x_i$ is exactly the
sample proportion of observed 1's; (b) states the assumptions under which this
estimator is appropriate; and (c) gives \emph{one} situation where it
misleads---for example, very small $N$, or all observed outcomes identical so
that $\hat{r}\in\{0,1\}$.

\subsection*{Going further (optional, ungraded)}

\noindent Place a $\mathrm{Beta}(a,b)$ prior on $r$ and find the MAP estimate.
Comment on how, for small $N$, the prior pulls the estimate away from the raw
sample proportion.

%%% Answer START %%%
\ans{
\subsection*{Solution}

\noindent\textbf{Assumptions and notation.} The $N$ samples are i.i.d.\ given
$r$, each $x_i \in \{0,1\}$, and $r \in (0,1)$. Here ``$\log$'' is the natural
logarithm and all sums run over $i = 1,\dots,N$.

\begin{enumerate}
    \item Because the samples are independent conditioned on $r$, the joint
        likelihood factorizes into the product of the per-sample masses:
        \[ L = \prod_{i = 1}^N r^{x_i}(1 - r)^{1 - x_i}. \]
    \item
        Maximizing $L$ is equivalent to maximizing $\log L$ (the logarithm is
        strictly increasing), so we work with the log-likelihood:
        \begin{align}
            \log L &= \sum_{i = 1}^N x_i \log r + (1-x_i) \log(1-r)\\
            \implies \frac{\partial}{\partial r}\log L &= \sum_{i = 1}^N \frac{x_i}{r} - \frac{1 - x_i}{1 - r}\\
            \implies \frac{\partial^2}{\partial r^2}\log L &= \sum_{i = 1}^N -\frac{x_i}{r^2}-\frac{1-x_i}{(1 - r)^2}
        \end{align}

        Every term of the second derivative is strictly negative for
        $r\in(0,1)$: since each $x_i\in\{0,1\}$, exactly one of the two
        fractions in each term is nonzero, and it enters with a negative sign.
        Hence $\log L$ is strictly concave on $(0,1)$, so it has at most one
        stationary point and any such point is its unique global maximum;
        therefore $L$ has a unique maximum.

    \item
        Setting the first derivative to zero at $r = \hat{r}$ and solving:
        \begin{align}
            \left.\frac{\partial \log L}{\partial r}\right|_{r = \hat{r}} &= 0\\
            \implies \sum_{i = 1}^N \frac{x_i}{\hat{r}} - \frac{1 - x_i}{1 - \hat{r}} &= 0\\
            \implies \frac{1}{\hat{r}}\sum_{i} x_i&=\frac{1}{1 - \hat{r}}\sum_i (1 - x_i)\\
            \implies (1 - \hat{r})\sum_i x_i &= \hat{r} \sum_i (1 - x_i)\\
            \implies \sum_i x_i - \hat{r} \sum_i x_i &= N\hat{r} - \hat{r}\sum_i x_i\\
            \implies \hat{r} &= \frac{\sum_i x_i}{N}
        \end{align}
\end{enumerate}

\medskip
\noindent\textbf{Interpretation.} The estimator $\hat{r}=\frac{1}{N}\sum_i x_i$
is the sample proportion of 1's because each $x_i\in\{0,1\}$, so $\sum_i x_i$
simply counts the observed successes and dividing by $N$ gives their fraction;
maximizing the Bernoulli likelihood reduces to matching $r$ to the observed
success rate. This is appropriate when the data really are i.i.d.\ draws from a
single fixed-$r$ Bernoulli. It misleads when $N$ is very small or when every
observation is identical: three heads in three tosses give $\hat{r}=1$, asserting
tails is impossible---an artifact of the tiny sample, not the truth. (A
$\mathrm{Beta}$ prior / MAP estimate, the optional extension, fixes exactly this
by pulling the estimate off the boundary.)

}
%%% Answer END %%%


\section{Workload}

How much time did you spend on this homework assignment outside of class?

\section*{Acknowledgments}

Cite all the people you've worked with on this homework, as well as any other
resources you used apart from the FCML textbook.

\end{document}
